
@@M0@@~;
montrons que @@M1@@ vérifie les conditions i) et ii).

i)~: @@M2@@.

ii)~: @@M3@@ d'après (1.6.1)~; or @@M4@@, et @@M5@@. Comme @@M6@@, et que @@M7@@ est un morphisme de foncteurs filtrés, on obtient ii).

Il faut enfin montrer que si @@M8@@ vérifie i) et ii), il existe @@M9@@ tel que @@M10@@.

On détermine d'abord la valeur de @@M11@@ pour l'élément @@M12@@ de @@M13@@ grâce à (1.6.1)~; on a donc par définition~: @@M14@@, les @@M15@@ étant connus par fonctorialité au moyen des morphismes @@M16@@.

On recolle ensuite les @@M17@@ pour @@M18@@ variable grâce au diagramme commutatif~:
@@M19@@
où @@M20@@ est déterminé par @@M21@@ si @@M22@@, @@M23@@, où @@M24@@ envoie @@M25@@ sur la @@M26@@-ième fonction symétrique des @@M27@@, et @@M28@@ envoie @@M29@@ sur la @@M30@@-ième fonction symétrique de @@M31@@ pour @@M32@@, et @@M33@@ sur @@M34@@.

En effet, on en déduit~:
@@M35@@
